% preamble-exam-zh.tex — 极简讲解页共享 preamble
% 目标：复刻「题目独立、提示轻框、路线箭头、主体双栏、答案轻收束」的讲义风格。

\documentclass{exam-zh}

% ---------- 基础配置 ----------
\examsetup{
  page/size = a4paper,
  page/show-head = false,
  page/show-foot = false,
  sealline/show = false,
  font = none,
  math-font = none,
}

% ---------- 包 ----------
\usepackage{geometry}
\geometry{top=10mm,bottom=14mm,left=12mm,right=10mm}
\AtBeginDocument{\newgeometry{top=10mm,bottom=14mm,left=12mm,right=10mm}}

\usepackage{xcolor}
\usepackage{needspace}
\usepackage{enumitem}
\usepackage{array}
\usepackage{tabularx}
\usepackage{tcolorbox}
\tcbuselibrary{skins,breakable}
\usepackage{paracol}

% ---------- 打印/屏幕颜色切换 ----------
% 用法：\documentclass[monochrome]{exam-zh} 可降级为灰度。
% 注意：不要使用 \if@xxx 放在 makeatother 外部；这里改成普通条件名。
\newif\ifeduprintcolor
\eduprintcolortrue
\makeatletter
\@ifclasswith{exam-zh}{monochrome}{\eduprintcolorfalse}{}
\makeatother

% ---------- 语义颜色 ----------
\ifeduprintcolor
  \definecolor{edu-blue}{RGB}{31,62,126}
  \definecolor{edu-blue-soft}{RGB}{232,238,250}
  \definecolor{edu-blue-bg}{RGB}{248,250,254}
  \definecolor{edu-ink}{RGB}{30,35,45}
  \definecolor{edu-muted}{RGB}{118,124,136}
  \definecolor{edu-line}{RGB}{209,218,235}
  \definecolor{edu-soft-bg}{RGB}{250,251,253}
  \definecolor{edu-red}{RGB}{168,58,58}
  \definecolor{edu-red-bg}{RGB}{255,247,245}
\else
  \definecolor{edu-blue}{gray}{0.15}
  \definecolor{edu-blue-soft}{gray}{0.90}
  \definecolor{edu-blue-bg}{gray}{0.98}
  \definecolor{edu-ink}{gray}{0.10}
  \definecolor{edu-muted}{gray}{0.45}
  \definecolor{edu-line}{gray}{0.82}
  \definecolor{edu-soft-bg}{gray}{0.97}
  \definecolor{edu-red}{gray}{0.28}
  \definecolor{edu-red-bg}{gray}{0.97}
\fi

% ---------- 全局排版 ----------
\setlength{\parindent}{0pt}
\setlength{\parskip}{0pt}
\linespread{1.16}
\renewcommand{\arraystretch}{1.15}

% ctex 通常提供 \kaishu；这里用它做教师旁白感。
\newcommand{\edukai}{\kaishu}

% ---------- 顶部元信息 ----------
\newcommand{\eduTopBar}[2]{%
  \noindent
  \begin{tabularx}{\linewidth}{@{}Xr@{}}
    {\color{edu-blue}\bfseries #1} &
    {\color{edu-blue}\bfseries #2}
  \end{tabularx}
  \vspace{8mm}
}

% ---------- 轻标题体系 ----------
% 只有真正的大结构才用它；不要给提示/路线滥用标题。
\newcommand{\eduSectionTitle}[1]{%
  \needspace{5\baselineskip}%
  \vspace{2mm}
  {\color{edu-blue}\bfseries\large #1}\par
  \vspace{3mm}
}

\newcommand{\eduSubTitle}[1]{%
  \needspace{4\baselineskip}%
  \vspace{1mm}
  {\color{edu-blue}\bfseries #1}\par
  \vspace{1mm}
}

% ---------- 小问题干：复用 exam (1)(2)(3) 格式 ----------
\newenvironment{eduSubQuestionItem}[1]{%
  \needspace{4\baselineskip}%
  \vspace{2mm}
  \par\noindent
  \hangindent=8mm\hangafter=1
  {\color{edu-blue}\bfseries #1}\enspace
  \begingroup
  \color{edu-ink}
}{%
  \endgroup
  \par
  \vspace{3mm}
}

% ---------- 辅助信息条（解题流程/关键想法/思考）楷体 ----------
\newenvironment{eduAuxItem}[1]{%
  \needspace{4\baselineskip}%
  \vspace{2mm}
  \par\noindent
  \hangindent=8mm\hangafter=1
  {\color{edu-blue}\edukai $\triangleright$~#1}\enspace
  \begingroup
  \color{edu-ink}
  \edukai
}{%
  \endgroup
  \par
  \vspace{4mm}
}

\newcommand{\eduColumnTitle}[2]{%
  {\color{edu-blue}\bfseries #1\hspace{1.5mm}#2}\par
  \vspace{2.5mm}
}

% ---------- 图标：不用外部 icon 字体，保证可移植 ----------
\newcommand{\eduIconBulb}{\(\circ\)}
\newcommand{\eduIconBook}{\(\square\)}
\newcommand{\eduIconCheck}{\(\checkmark\)}
\newcommand{\eduIconPin}{\(\diamond\)}
\newcommand{\eduIconNote}{\(\triangleright\)}

% ---------- 题目区 ----------
% 题目区是页面入口：弱标签 + 一条长蓝线 + 大题干。
\newenvironment{eduProblemBlock}[1][题目]{%
  \needspace{8\baselineskip}
  {\small\color{edu-muted}#1}\par
  \vspace{1.5mm}
  {\color{edu-blue}\hrule height 0.55pt}
  \vspace{4.5mm}
  \begingroup
  \color{edu-ink}
  \large
  \setlength{\parskip}{2.2mm}
  \linespread{1.20}\selectfont
}{%
  \par
  \endgroup
  \vspace{6mm}
}

\newenvironment{eduSubQuestions}{%
  \begin{enumerate}[
    label=(\arabic*),
    leftmargin=8mm,
    itemsep=2.0mm,
    topsep=2.0mm,
    parsep=0pt
  ]
}{%
  \end{enumerate}
}

% ---------- 读题提示：轻框，不做同级标题 ----------
\newtcolorbox{eduReadTipBox}{
  enhanced,
  breakable,
  colback=white,
  colframe=edu-blue!45,
  boxrule=0.45pt,
  arc=1.6mm,
  left=10mm,
  right=4mm,
  top=5mm,
  bottom=5mm,
  before skip=1mm,
  after skip=7mm,
  fontupper=\edukai\normalsize,
  colupper=edu-ink,
  overlay={
    \node[anchor=north west, text=edu-blue] at ([xshift=3.2mm,yshift=-3.2mm]frame.north west)
    {\small\eduIconNote};
  }
}

\newenvironment{eduTipList}{%
  \begin{itemize}[
    label=\textbullet,
    leftmargin=5mm,
    itemsep=1.2mm,
    topsep=0pt,
    parsep=0pt
  ]
}{%
  \end{itemize}
}

% ---------- 解题路线：虚线框 + 文字箭头 ----------
\newenvironment{eduRouteFlow}{%
  \needspace{4\baselineskip}%
  \vspace{1mm}
  \begin{center}
  \normalsize
}{%
  \end{center}
  \vspace{7mm}
}
\newcommand{\eduRouteStep}[1]{%
  {\color{edu-ink}#1}%
}
\newcommand{\eduRouteArrow}{%
  \;$\boldsymbol{\to}$\;%
}

% ---------- 主体讲解：使用 paracol 实现跨页自适应双栏 ----------
\newenvironment{eduDualExplain}{%
  \needspace{6\baselineskip}% 防止标题孤行
  \columnratio{0.60, 0.40}%
  \setlength{\columnsep}{11mm}%
  \columnseprule=0.45pt%
  \colseprulecolor{edu-blue}%
  \begin{paracol}{2}% 开启双栏
}{%
  \end{paracol}% 结束双栏
  \vspace{5mm}
}

\newenvironment{eduExplainLeft}{%
  \normalsize\color{edu-ink}%
}{%
  \switchcolumn
}

\newcommand{\eduColumnDivider}{}

\newenvironment{eduExplainRight}{%
  \normalsize\color{edu-ink}%
}{%
}

\newenvironment{eduNumberList}{%
  \begin{enumerate}[
    label=\arabic*.,
    leftmargin=6mm,
    itemsep=4.8mm,
    topsep=0pt,
    parsep=0pt
  ]
}{%
  \end{enumerate}
}

\newenvironment{eduBulletList}{%
  \begin{itemize}[
    label=\textbullet,
    leftmargin=5mm,
    itemsep=1.2mm,
    topsep=0pt,
    parsep=0pt
  ]
}{%
  \end{itemize}
}

% ---------- 底部双栏：答案 + 方法提醒 ----------
\newenvironment{eduBottomGrid}{%
  \columnratio{0.32, 0.68}%
  \setlength{\columnsep}{11mm}%
  \columnseprule=0pt%
  \begin{paracol}{2}%
}{%
  \end{paracol}%
}

\newenvironment{eduSummaryLeft}{%
}{%
}

\newcommand{\eduSummaryDivider}{%
  \switchcolumn
}

\newenvironment{eduSummaryRight}{%
}{%
}

% ---------- 答案/方法提醒：轻量收束 ----------
\newtcolorbox{eduSummaryBlock}[2]{
  enhanced,
  breakable,
  colback=edu-blue-bg,
  colframe=edu-line,
  boxrule=0.45pt,
  arc=1.2mm,
  left=3.5mm,
  right=3.5mm,
  top=3mm,
  bottom=3mm,
  before skip=1mm,
  after skip=2.5mm,
  title={\color{edu-blue}\bfseries #1\hspace{1.5mm}#2},
  coltitle=edu-blue,
  attach title to upper={\par\vspace{1.5mm}},
  fontupper=\normalsize
}

% ---------- 例题单栏解答卡片 ----------
\newtcolorbox{eduSolutionBox}[1]{
  enhanced,
  breakable,
  colback=white,
  colframe=edu-blue!40,
  boxrule=0.55pt,
  arc=1.2mm,
  left=4mm,
  right=4mm,
  top=3.5mm,
  bottom=3.5mm,
  before skip=2mm,
  after skip=3mm,
  title={\color{edu-blue}\bfseries \eduIconBook\hspace{1.5mm}#1},
  coltitle=edu-blue,
  attach title to upper={\par\vspace{1.5mm}},
  fontupper=\normalsize
}

\newenvironment{eduPlainList}{%
  \begin{enumerate}[
    label=(\arabic*),
    leftmargin=7mm,
    itemsep=1.2mm,
    topsep=0pt,
    parsep=0pt
  ]
}{%
  \end{enumerate}
}

% ---------- 兼容旧 block 的轻量盒子 ----------
\newtcolorbox{eduSoftBox}{
  enhanced,
  breakable,
  colback=edu-soft-bg,
  colframe=edu-line,
  boxrule=0.35pt,
  arc=1mm,
  left=3mm,
  right=3mm,
  top=2mm,
  bottom=2mm,
  before skip=1mm,
  after skip=2mm,
  fontupper=\small
}

\newtcolorbox{eduMistakeBox}{
  enhanced,
  breakable,
  colback=edu-red-bg,
  colframe=edu-red!40,
  boxrule=0.35pt,
  arc=1mm,
  left=3mm,
  right=3mm,
  top=2.5mm,
  bottom=2.5mm,
  before skip=1mm,
  after skip=2mm,
  fontupper=\normalsize
}

\newtcolorbox{eduLegacySolution}{
  enhanced,
  breakable,
  colback=edu-soft-bg,
  colframe=edu-line,
  boxrule=0.35pt,
  arc=1mm,
  left=3mm,
  right=3mm,
  top=2mm,
  bottom=2mm,
  before skip=-2mm,
  after skip=4mm,
  fontupper=\small
}

\newtcolorbox{eduTeacherNote}{
  enhanced,
  breakable,
  colback=edu-soft-bg,
  colframe=edu-line,
  boxrule=0pt,
  borderline west={1.5pt}{0pt}{edu-muted},
  left=2.5mm,
  right=2.5mm,
  top=2mm,
  bottom=2mm,
  before skip=1mm,
  after skip=2mm,
  fontupper=\footnotesize,
  colupper=edu-muted
}

% ---------- 答题区 ----------
\newcommand{\answerarea}[1][40mm]{%
  \vspace{3mm}%
  \begin{tcolorbox}[
    enhanced,
    breakable,
    colback=white,
    colframe=edu-line,
    boxrule=0.55pt,
    arc=0.8mm,
    height=#1,
    valign=top,
  ]
  \end{tcolorbox}%
}

\examsetup{
  question/show-answer = true,
  fillin/show-answer = true,
  solution/show-solution = show-stay,
}

\begin{document}


\begin{eduProblemBlock}[题目]
已知大圆半径为 $R$，小圆半径为 $r$（即 $R \ge r$），两圆圆心距为 $d$。
\begin{enumerate}[label=(\arabic*)]
  \item \textbf{【已知半径与圆心距求位置关系】} 已知 $\odot O_1$ 的半径为 $3\text{cm}$，$\odot O_2$ 的半径为 $4\text{cm}$。
  \begin{enumerate}[label=\arabic*.]
    \item 若圆心距 $O_1O_2 = 8\text{cm}$，求两圆的位置关系；
    \item 若圆心距 $O_1O_2 = 5\text{cm}$，求两圆的位置关系。
  \end{enumerate}
  \item \textbf{【已知半径与位置关系求圆心距】} 已知两圆的半径分别为 $2$ 和 $5$，若两圆\textbf{相切}，求两圆的圆心距 $d$。
  \item \textbf{【已知单半径、圆心距与位置关系求另一半径】} 已知 $\odot A$ 的半径为 $4$，$\odot A$ 与 $\odot B$ \textbf{相切}，两圆的圆心距为 $3$，求 $\odot B$ 的半径 $r$。
\end{enumerate}

\end{eduProblemBlock}









\begin{eduAuxItem}{思路导航}
计算关键边界值：和值 $R+r$ 与差值 $R-r$$\;\to\;$对比圆心距 $d$ 与边界值，确定几何位置关系$\;\to\;$相切讨论：外切代入 $d = R+r$，内切代入 $d = |R-r|$$\;\to\;$内切大圆讨论：分 $r_A > r_B$ 或 $r_B > r_A$ 讨论$\;\to\;$排除非法解（半径必须大于 0）并汇总答案\end{eduAuxItem}




\begin{eduAuxItem}{核心思路}
\textbf{数轴边界模型（位置关系区间化）}：
将圆心距 $d$ 映射在数轴上，两半径的“和” $R+r$ 与“差” $R-r$ （设 $R \ge r$）即为核心边界点。
- 边界点 1：$d = R + r \iff$ 外切
- 边界点 2：$d = R - r \iff$ 内切
当题目涉及“相切”时，须立即对应 $d = R+r$ 或 $d = R-r$ 两种情况进行分类讨论。
其他关系为区间：
- “相交”对应开区间：$R-r < d < R+r$
- “内含”对应区间：$0 \le d < R-r$
\end{eduAuxItem}




\begin{eduSubQuestionItem}{(1)}
已知 $\odot O_1$ 的半径为 $3\text{cm}$，$\odot O_2$ 的半径为 $4\text{cm}$。
若圆心距 $O_1O_2 = 8\text{cm}$，求两圆的位置关系；若圆心距 $O_1O_2 = 5\text{cm}$，求两圆的位置关系。
\end{eduSubQuestionItem}

\begin{eduSolutionBox}{解答}
  \begingroup
  \setlength{\parskip}{2.5mm}
由题意，设 $R = 4\text{cm}, r = 3\text{cm}$，则 $R+r = 7\text{cm}$，$R-r = 1\text{cm}$。\par
① 当圆心距 $d = 8\text{cm}$ 时，因 $d > R+r$，故两圆\textbf{外离}。\par
② 当圆心距 $d = 5\text{cm}$ 时，因 $R-r < d < R+r$，故两圆\textbf{相交}。\par
  \endgroup
  \vspace{2.5mm}
  {\color{edu-blue}\bfseries\small 核心}\par\vspace{1.5mm}
  \begin{eduBulletList}
    \item 求解前应先求出两半径之“和”与“差”，作为判定两圆位置关系的基准线。  \end{eduBulletList}
\end{eduSolutionBox}




\begin{eduSubQuestionItem}{(2)}
已知两圆的半径分别为 $2$ 和 $5$，若两圆相切，求两圆的圆心距 $d$。\end{eduSubQuestionItem}

\begin{eduSolutionBox}{解答}
  \begingroup
  \setlength{\parskip}{2.5mm}
因两圆相切，分外切和内切两种情况讨论：\par
① 两圆外切时，圆心距等于半径之和：$d = R+r = 5+2=7$；\par
② 两圆内切时，圆心距等于半径之差：$d = R-r = 5-2=3$。\par
综上，两圆的圆心距 $d$ 为 $3$ 或 $7$。\par
  \endgroup
  \vspace{2.5mm}
  {\color{edu-blue}\bfseries\small 关键}\par\vspace{1.5mm}
  \begin{eduBulletList}
    \item 在未指明相切方式时，“相切”必须包含“外切”和“内切”两种情形，避免遗漏。  \end{eduBulletList}
\end{eduSolutionBox}




\begin{eduSubQuestionItem}{(3)}
已知 $\odot A$ 的半径为 $4$，$\odot A$ 与 $\odot B$ 相切，两圆的圆心距为 $3$，求 $\odot B$ 的半径 $r$。\end{eduSubQuestionItem}

\begin{eduSolutionBox}{解答}
  \begingroup
  \setlength{\parskip}{2.5mm}
设 $\odot B$ 的半径为 $r(r>0)$。因 $\odot A$ 与 $\odot B$ 相切，分以下两种情况讨论：\par
① 两圆外切时，圆心距 $d = r_A + r \implies 3 = 4 + r \implies r = -1 < 0$（不合题意，舍去）；\par
② 两圆内切时，圆心距 $d = |r_A - r|$。分情况讨论大圆归属：\par
\quad 情况 a（$\odot A$ 为大圆，即 $4 > r$）：$d = r_A - r \implies 3 = 4 - r \implies r = 1$；\par
\quad 情况 b（$\odot B$ 为大圆，即 $r > 4$）：$d = r - r_A \implies 3 = r - 4 \implies r = 7$。\par
综上，$\odot B$ 的半径为 $1$ 或 $7$。\par
  \endgroup
  \vspace{2.5mm}
  {\color{edu-blue}\bfseries\small 警示}\par\vspace{1.5mm}
  \begin{eduBulletList}
    \item 若大圆身份未明，内切仍有两解。这是填空题中的易错考点，必须对大小圆的归属进行完整讨论。  \end{eduBulletList}
\end{eduSolutionBox}




\begin{eduBottomGrid}
  \begin{eduSummaryLeft}
    \begin{eduSummaryBlock}{\eduIconCheck}{通关答案汇总}
    \begin{eduPlainList}
      \item (1) $O_1O_2=8\text{cm}$ 时外离； $O_1O_2=5\text{cm}$ 时相交。      \item (2) 圆心距 $d$ 为 $3$ 或 $7$。      \item (3) $\odot B$ 的半径为 $1$ 或 $7$。    \end{eduPlainList}
    \end{eduSummaryBlock}
  \end{eduSummaryLeft}
  \eduSummaryDivider
  \begin{eduSummaryRight}
    \begin{eduSummaryBlock}{\eduIconPin}{“知二推一”核心口诀}
    \begin{eduBulletList}
      \item 两圆相切定边界，外切内切分两界。      \item 已知半径求半径，大圆小圆要警惕。      \item 外切负解须舍去，内切讨论防漏解。    \end{eduBulletList}
    \end{eduSummaryBlock}
  \end{eduSummaryRight}
\end{eduBottomGrid}




\begin{eduMistakeBox}
\eduSubTitle{漏掉内切情况}\textbf{漏掉内切情况}：
“相切”包括“外切”和“内切”两种情形，极易遗漏内切。
防坑策略：若未明确指明相切类型，看见“相切”应立刻列出外切（$d=R+r$）和内切（$d=R-r$）两种方程。
\end{eduMistakeBox}




\begin{eduMistakeBox}
\eduSubTitle{忽略负半径必须舍去}\textbf{忽略负半径舍去}：
易将计算出的负值半径写进答案。
防坑策略：在几何学中，圆半径必须为正数（$r > 0$）。若算出负数解，须作“不合题意，舍去”处理。
\end{eduMistakeBox}




\begin{eduMistakeBox}
\eduSubTitle{漏掉内切中的大圆分类}\textbf{漏掉内切中的大圆归属讨论}：
默认已知圆为大圆，遗漏新求圆为大圆的情况。
防坑策略：在内切且未知大圆归属时，必须列出 $d = R_{已知} - r_{未知}$ 和 $d = r_{未知} - R_{已知}$ 两种可能进行完整求解。
\end{eduMistakeBox}




\begin{eduAuxItem}{思路启发}
若将例题 3 中的圆心距改为 5，求 $\odot B$ 的半径的有效解个数。（提示：列方程计算，检查外切解是否为正数）\end{eduAuxItem}




\begin{eduAuxItem}{思路启发}
若两圆位置关系为非同心圆的“内含”，圆心距 $d$ 的取值范围如何表示？（提示：注意 $d > 0$ 的约束）\end{eduAuxItem}




\end{document}
